% Cote 119: offprints, reports, courses, CV and letters (typed and printed part): archive cote 119, site page 179 (PDF page 180).
% typescript. Draft compared: none.
% AI-transcribed by Claude Opus 5.5 (claude-opus-5-5, adaptive thinking, effort high), Message Batches API, batch r119_full (msgbatch_01DwZNm5Y7YsKBWyAN4sFs8w), 2026-11-13; not yet checked by eye.
\scanpage{119/179}
\mspage{- 18 -}

On se reportera avec profit à la feuille d'exercices n\textsuperscript{o} 4. II 2\textsuperscript{o} b \unsure{(}et 3\textsuperscript{o} a).

On peut donc en ajoutant des constantes aux $f_j$ (c'est-à-dire en
modifiant une fois de plus les $F_{i,j}$ choisis) rendre $f_j(t,u)$ et $f_{j+1}(t,u)$
égales pour $u = u_{j+1}$. La fonction $f(t,u)$ définie par $f(t,u) = f_j(t,u)$
pour $u \in [u_j, u_{j+1}]$ est bien continue et vérifie $F_{i,j}(\delta(t,u)) = f(t,u)$
pour $t_i \leq t \leq t_{i+1}$, $u_j \leq u \leq u_{j+1}$ ce qui démontre l'existence de f.
Si $f_1$ et $f_2$ sont deux telles fonctions, au voisinage de tout point $(t_0,u_0)$
du rectangle $[a,b] \times [a',b']$ on a
$f_1(t,u) - f_2(t,u) = F_1(\delta(t,u)) - F_2(\delta(t,u))$ qui est constant. La
fonction $f_1-f_2$ \unsure{est} donc ``localement\unsure{''} constante''. Comme le rectangle
$[a,b] \times [a',b']$ est connexe on voit facilement que $f_1-f_2$ est constante
d'où la dernière assertion.

Le cas où pour chaque $t$ la fonction $u \to \delta(t,u)$ est constante nous
fournit le résultat intéressant :

6.6. \uline{Corollaire}

Soit $\gamma : [a,b \to U$ un \unsure{clemin|chemin} contenu dans l'ouvert U ($\gamma$ est une
application continue) et soit $\omega$ une forme différentielle fermée dans U. Il
existe une fonction f continue sur $[a,b]$ telle que au voisinage de tout
point $\gamma\ (t_0)$ ($a \leq t_0 \leq b$) il existe une primitive F de $\omega$ telle que
$F(\gamma(t)) = f(t)$ . f est unique à l'addition d'une constante près.
(On dira que f est une primitive de $\omega$ le long de $\gamma$).

6.7. \uline{Remarque}.

Si $\gamma$ est différentiable par morceaux, f une primitive de $\omega$
suivant $\gamma$ on vérifie facilement que $\int_\gamma \omega = f(b) - f(a)$.
D'après (6.6) cette formule donne une définition de $\int_\gamma \omega$ dans le cas où $\gamma$

\margin{179}
